\section*{الفصل \ref{ch:g11:subduction-zones} --- نطاقات الاندساس}

\begin{solution}{exo:g11:subduction-zones:1}
\omterm{def:g11:subduction-zones:subduction}{الاندساس} هو غوص صفيحة من \omterm{def:g11:ocean-ridges:oceanic-lithosphere}{الغلاف الصخري المحيطي} تحت صفيحة
أخرى، عائدةً إلى الوشاح؛ \omterm{def:g11:subduction-zones:subduction}{ونطاق الاندساس} حد متقارب بين صفيحتين
يحدث فيه \omterm{def:g11:subduction-zones:subduction}{الاندساس}، ويَسِمه خندق. \omterm{def:g11:ocean-ridges:oceanic-lithosphere}{والغلاف الصخري المحيطي} هو الذي يغوص: فحين
يقدم، يصير باردًا وسميكًا وأكثف من \omterm{def:g11:lithosphere-asthenosphere:lithosphere}{الغلاف اللدن} الذي تحته. أما
القشرة القارية فأخف من أن تغوص.
\end{solution}

\begin{solution}{exo:g11:subduction-zones:2}
(1) على سبيل المثال منطقة تونغا--فيجي، وغرب المحيط الهادئ من اليابان إلى
جزر ماريانا، وإندونيسيا، وأمريكا الجنوبية. (2) دائمًا على جانب
الصفيحة الراكبة، ولا تقع أبدًا على الجانب المحيطي. (3) ضحلة: فكل زلازل
ظهرة وسط الأطلسي على الخريطة أقل عمقًا من \qty{70}{km}.
\end{solution}

\begin{solution}{exo:g11:subduction-zones:3}
$7031 \times 0.21 \approx 1480$ زلزالًا أعمق من \qty{70}{km}؛
و$7031 \times 0.069 \approx 485$ أعمق من \qty{300}{km}.
\end{solution}

\begin{solution}{exo:g11:subduction-zones:4}
\omterm{def:g11:subduction-zones:eclogite}{الشيست الأزرق} صخر متحول مصنوع من بازلت \omterm{def:g11:ocean-ridges:oceanic-lithosphere}{القشرة المحيطية}
المدفون تحت ضغط مرتفع مع بقائه باردًا نسبيًا، ويلوّنه بالأزرق
الغلوكوفان؛ \omterm{def:g11:subduction-zones:eclogite}{والإكلوجيت} هو الصخر المتحول الكثيف الذي يصير إليه البازلت \omterm{def:g11:ocean-ridges:gabbro}{والغابرو}
على أعماق الوشاح، وهو مصنوع من العقيق الأحمر والبيروكسين الأخضر. وكلاهما آتٍ من
\omterm{def:g11:ocean-ridges:oceanic-lithosphere}{القشرة المحيطية} ويحتفظ بتركيبها الكيميائي؛ أما ما يختلف فهو
المعادن، التي تكوّنت في الحالة الصلبة تحت الضغط ودرجة الحرارة الجديدين،
والكثافة.
\end{solution}

\begin{solution}{exo:g11:subduction-zones:5}
قاع المحيط: \omterm{def:g11:lithosphere-asthenosphere:lithosphere}{غلاف صخري} قديم يتحرك نحو الخندق. الخندق: تنثني الصفيحة
نحو الأسفل وتبدأ بالغوص. مقدمة القوس: تتراكم الرواسب المكشوطة عن الصفيحة الغائصة؛
وتحتها تنزلق الصفيحتان إحداهما على الأخرى على امتداد الحد بين الصفيحتين،
حيث تحدث الزلازل الكبرى. القوس البركاني: خط من البراكين فوق
\omterm{def:g11:subduction-zones:slab}{اللوح المندس}. ما وراء القوس: خلف بعض الأقواس، تتمطط الصفيحة الراكبة.
\end{solution}

\begin{solution}{exo:g11:subduction-zones:6}
$\tan\delta = (508 - 135)/(800 - 300) = 373/500 \approx 0.746$، إذن
$\delta \approx \ang{37}$.
\end{solution}

\begin{solution}{exo:g11:subduction-zones:7}
(1) $\tan\ang{45} = 1$: على عمق \qty{300}{km}. (2) على بعد \qty{670}{km} من
الخندق. (3) $670\sqrt{2} \approx \qty{948}{km}$ من \omterm{def:g11:subduction-zones:slab}{اللوح المندس}.
\end{solution}

\begin{solution}{exo:g11:subduction-zones:8}
(1) $(74.6 + 70.0 + 65.5 + 63.1 + 52.6)/5 \approx \qty{65.2}{mm}$ في السنة؛
و$(47.6 + 31.0 + 29.2 + 22.8 + 19.8 + 15.2 + 14.9 + 10.6)/8 \approx
\qty{23.9}{mm}$ في السنة؛ والنسبة $65.2/23.9 \approx 2.7$. (2) تتحرك الصفائح
أسرع ما يكون حين يتدلى منها \omterm{def:g11:subduction-zones:slab}{لوح مندس}: \omterm{def:g11:subduction-zones:slab-pull}{فشد اللوح المندس} هو
أقوى القوى التي تحرّك الصفائح، كما استنتج فورسايث وأويدا.
\end{solution}

\begin{solution}{exo:g11:subduction-zones:9}
(1) \qty{83}{mm} تساوي \qty{0.083}{m}: $50/0.083 \approx 600$ سنة
و$60/0.083 \approx 720$ سنة من حركة الصفيحة. (2) من 869 إلى 2011 تمر 1142
سنة: من الرتبة نفسها. فإذا بقي الحد بين الصفيحتين مقفلًا بين الزلازل
الكبرى، فإن انزلاق 2011 أطلق عدة قرون من الحركة المختزنة.
\end{solution}

\begin{solution}{exo:g11:subduction-zones:10}
(1) اليابان: $6400 - 3200\,e^{-133/62.8} \approx \qty{6015}{m}$؛ تونغا:
$6400 - 3200\,e^{-99/62.8} \approx \qty{5740}{m}$؛ الأنديز:
$2500 + 350\sqrt{49} = \qty{4950}{m}$. (2) القاع الذي يبلغ اليابان: فهو
الأقدم، ومن ثم \omterm{def:g11:lithosphere-asthenosphere:lithosphere}{الغلاف الصخري} الأبرد والأكثف، والقاع الأعمق.
\end{solution}

\begin{solution}{exo:g11:subduction-zones:11}
(1) $3380 \times 1.06 \approx \qty{3583}{kg/m^3}$. (2)
$(3583 - 2900)/2900 \approx 0.24$: أي ارتفاع بنحو \qty{24}{\percent}.
(3) قبل التحول تكون القشرة أخف من الوشاح بمقدار
$3380 - 2900 = \qty{480}{kg/m^3}$؛ وبعده تكون أكثف منه بنحو
\qty{203}{kg/m^3}. فلم تعد تسند \omterm{def:g11:subduction-zones:slab}{اللوح المندس}: بل تشده نحو الأسفل.
\end{solution}

\begin{solution}{exo:g11:subduction-zones:12}
(1) $t = 6 \times 10^6 \times 3.16 \times 10^7 \approx 1.9 \times 10^{14}$~\unit{s}؛
$\sqrt{10^{-6} \times 1.9 \times 10^{14}} \approx \qty{1.4e4}{m}$، أي نحو
\qty{14}{km}. (2) أقل بكثير من \qty{100}{km}: ففي 6 ملايين سنة لا يسخن إلا
الغلاف الخارجي الرقيق من \omterm{def:g11:subduction-zones:slab}{اللوح المندس}؛ ويبقى داخله باردًا. (3) الصخر البارد
صلب ويمكن أن ينكسر، بينما ينساب الوشاح الساخن المحيط به؛ والزلازل
العميقة تحدث داخل \omterm{def:g11:subduction-zones:slab}{اللوح المندس} البارد.
\end{solution}

\begin{solution}{exo:g11:subduction-zones:13}
(1) $\sqrt{1017^2 + 670^2} \approx \qty{1218}{km}$. (2) \qty{64}{mm} في السنة
تساوي \qty{64}{km} في كل مليون سنة: $1218/64 \approx 19$ مليون سنة.
(3) أطول بنحو ثلاث مرات منها تحت تونغا: \omterm{def:g11:subduction-zones:slab}{فاللوح المندس} تحت الأنديز
يغوص أبطأ وبميل ألطف.
\end{solution}

\begin{solution}{exo:g11:subduction-zones:14}
(1) شريحة عرضها \qty{1}{m} على امتداد الخندق حجمها
$10^5 \times 10^6 \times 1 = \qty{1e11}{m^3}$؛ وفائض كتلتها
$36 \times 10^{11} = \qty{3.6e12}{kg}$. (2) فائض وزنها
$3.6 \times 10^{12} \times 9.8 \approx \qty{3.5e13}{N}$ لكل متر من
الخندق. (3) \omterm{def:g11:subduction-zones:slab}{اللوح المندس} متصل ببقية الصفيحة، وهي صلبة:
فتنتقل القوة إلى الصفيحة كلها، فتنجرّ نحو
الخندق، لا سيما أن الخنادق تمتد آلاف الكيلومترات.
\end{solution}

\begin{solution}{exo:g11:subduction-zones:15}
(1) $(23.6 - 16.5) \times 111 \approx \qty{790}{km}$. (2) يزداد المعدل بمقدار
$237 - 104 = \qty{133}{mm}$ في السنة على امتداد \qty{790}{km}، أي نحو \qty{17}{mm} في
السنة لكل مئة كيلومتر. (3) التقارب هو سرعة الصفيحة
الهادئة مطروحًا منها سرعة القوس. ولأن سرعة الصفيحة الهادئة تكاد
تكون واحدة عند النقطتين، فلا بد أن القوس نفسه يتحرك نحو
الصفيحة الهادئة، أسرع فأسرع نحو الشمال: فمنطقة ما وراء القوس التي خلفه
تتمطط، في الشمال أكثر منها في الجنوب.
\end{solution}

\begin{solution}{pb:g11:subduction-zones:1}
\textbf{1.} نحو \qty{671}{km}.

\textbf{2.} على الجانب المحيطي لم تغص الصفيحة بعد: \omterm{def:g11:lithosphere-asthenosphere:lithosphere}{فالغلاف الصخري}
هناك قريب من السطح. والزلازل العميقة تحدث داخل \omterm{def:g11:subduction-zones:slab}{اللوح المندس}، الذي
لا يقع إلا تحت الصفيحة الراكبة.

\textbf{3.} $z(300) = -24 + 273 = \qty{249}{km}$
و$z(700) = -24 + 637 = \qty{613}{km}$؛ $\tan\delta = 364/400 = 0.91$، إذن
$\delta \approx \ang{42}$.

\textbf{4.} $x = (100 + 24)/0.91 \approx \qty{136}{km}$
و$x = (400 + 24)/0.91 \approx \qty{466}{km}$.

\textbf{5.} لو كان خطًا مستقيمًا بهذا الميل يبدأ من الخندق لمرّ
بالنقطة $z = 0$ عند $x = 0$؛ لكنه يقع فوقها. \omterm{def:g11:subduction-zones:slab}{فاللوح المندس} ليس مستقيمًا: فهو قرب
الخندق يميل ميلًا ألطف، ثم يشتد ميله في العمق.

\textbf{6.} \qty{16.8}{cm} في السنة، و\qty{168}{km} في كل مليون سنة.

\textbf{7.} $168/69 \approx 2.4$ مرة أسرع.

\textbf{8.} التقارب عبر الخندق هو سرعة الصفيحة الهادئة
مطروحًا منها سرعة القوس على الجانب الآخر. فإذا تحرك القوس نحو
الصفيحة الهادئة، كما يفعل بينما تتمطط منطقة ما وراء القوس التي خلفه، فإن
السرعتين تتجمعان، ويكون التقارب أكبر من سرعة
الصفيحة الهادئة وحدها.

\textbf{9.} $6400 - 3200\,e^{-99/62.8} \approx \qty{5740}{m}$.

\textbf{10.} بعد نحو 60 مليون سنة يصير \omterm{def:g11:ocean-ridges:oceanic-lithosphere}{الغلاف الصخري المحيطي} أكثف، في
مجمله، من \omterm{def:g11:lithosphere-asthenosphere:lithosphere}{الغلاف اللدن} الذي تحته؛ وعند 99 مليون سنة يكون أكثف
أيضًا. فما إن ينزل طرفه حتى يسحب وزنه البقية إلى الداخل.

\textbf{11.} $(670 + 24)/0.91 \approx \qty{763}{km}$.

\textbf{12.} $\sqrt{763^2 + 670^2} \approx \qty{1015}{km}$.

\textbf{13.} $1015/168 \approx 6.0$ مليون سنة.

\textbf{14.} $1015/69 \approx 14.7$ مليون سنة.

\textbf{15.} $3380 \times 1.06 \approx \qty{3583}{kg/m^3}$: فالقشرة، التي كانت
أخف من الوشاح، تصير أكثف منه، فيشد \omterm{def:g11:subduction-zones:slab}{اللوح} كله
بقوة أكبر.

\textbf{16.} $t \approx 6.0 \times 10^6 \times 3.16 \times 10^7 \approx
1.9 \times 10^{14}$~\unit{s}، و$\sqrt{\kappa t} \approx \qty{14}{km}$.

\textbf{17.} \qty{14}{km} قليلة إلى جانب \qty{100}{km}: فلم يسخن إلا الجزء الخارجي
من \omterm{def:g11:subduction-zones:slab}{اللوح المندس}؛ ويبقى داخله باردًا حين يبلغ
\qty{670}{km}.

\textbf{18.} يبقى \omterm{def:g11:subduction-zones:slab}{اللوح المندس} البارد صلبًا بما يكفي لينكسر، حتى طرفه،
بينما ينساب الوشاح الساخن المحيط به: فالزلازل تحت تونغا تحدث
داخل \omterm{def:g11:subduction-zones:slab}{اللوح المندس}، حتى عمق نحو \qty{671}{km}.

\textbf{19.} تُظهر الموجات الزلزالية تغيرًا حادًا في الوشاح على عمق نحو
\qty{670}{km}: إذ تقفز سرعاتها هناك.

\textbf{20.} نحو \textbf{6 ملايين سنة} من الخندق حتى
\qty{670}{km}.
\end{solution}
